\label{ch:resultados}
En este capítulo se presentarán y analizarán los resultados obtenidos a partir de la experimentación y evaluación comparativa de los distintos modelos de lenguaje evaluados sobre el dataset de referencia. 

Se detallará el rendimiento de clasificación binaria (coincidencia exitosa frente a discrepancias) mediante matrices de confusión, calculando de manera desagregada la precisión, sensibilidad y el valor F1-Score para cada modelo evaluado. Asimismo, se expondrán las puntuaciones de calidad otorgadas a las justificaciones administrativas generadas, evaluadas mediante la metodología \textit{LLM-as-a-judge} en una escala de aceptabilidad semántica.

Finalmente, se presentará el análisis comparativo de la eficiencia operativa en producción, correlacionando los tiempos de respuesta promedio por consulta con el costo financiero por lote de procesamiento. Este análisis cuantitativo servirá de sustento para justificar objetivamente la selección del modelo de lenguaje óptimo que equilibra la alta precisión administrativa con la viabilidad económica de la Clínica Dental Docente de la UPCH.
